% ═══════════════════════════════════════════════════════════════
% AB-04: Energieumwandlung & Erhaltungssatz
% Physik Klasse 10 MR
% ═══════════════════════════════════════════════════════════════

\documentclass[11pt, a4paper]{article}

\usepackage[utf8]{inputenc}
\usepackage[T1]{fontenc}
\usepackage[ngerman]{babel}

\usepackage{helvet}
\renewcommand{\familydefault}{\sfdefault}

\usepackage{microtype}
\usepackage[margin=1.5cm, top=2cm, bottom=1.5cm]{geometry}
\usepackage{enumitem}
\usepackage{tabularx}
\usepackage{booktabs}
\usepackage{array}
\usepackage{tikz}
\usetikzlibrary{calc, arrows.meta, shapes.geometric, positioning}

\usepackage{amsmath, amssymb}
\usepackage{siunitx}
\sisetup{locale = DE, per-mode = symbol}

\usepackage{fancyhdr}
\usepackage{lastpage}
\usepackage{needspace}

\usepackage{xcolor}
\usepackage{tcolorbox}
\tcbuselibrary{skins}

% ═══════════════════════════════════════════════════════════════
% FARBEN
% ═══════════════════════════════════════════════════════════════

\definecolor{physik}{HTML}{2c5aa0}
\definecolor{hellgrau}{HTML}{F5F5F5}
\definecolor{mittelgrau}{HTML}{888888}
\definecolor{rahmen}{HTML}{CCCCCC}
\definecolor{Epot}{HTML}{2c5aa0}
\definecolor{Ekin}{HTML}{c62828}
\definecolor{Etherm}{HTML}{888888}
\definecolor{bnefarbe}{HTML}{2a7a4b}

\colorlet{fachfarbe}{physik}

% ═══════════════════════════════════════════════════════════════
% VARIABLEN
% ═══════════════════════════════════════════════════════════════

\newcommand{\fach}{Physik}
\newcommand{\thema}{Energieumwandlung \& Erhaltungssatz}
\newcommand{\klasse}{10 MR}
\newcommand{\stunde}{04}
\newcommand{\maxpunkte}{30}

% ═══════════════════════════════════════════════════════════════
% KOPF- UND FUSSZEILE
% ═══════════════════════════════════════════════════════════════

\pagestyle{fancy}
\fancyhf{}
\renewcommand{\headrulewidth}{0.4pt}
\renewcommand{\footrulewidth}{0pt}
\lhead{\small \fach{} Klasse \klasse{} | \thema}
\rhead{\small Name: \underline{\hspace{3.5cm}}}
\cfoot{\small\textcolor{mittelgrau}{Seite \thepage{} von \pageref{LastPage}}}

% ═══════════════════════════════════════════════════════════════
% BOXEN
% ═══════════════════════════════════════════════════════════════

\newtcolorbox{merkkasten}[1][]{
    colback=hellgrau,
    colframe=hellgrau,
    coltitle=black,
    colbacktitle=hellgrau,
    boxrule=0pt,
    arc=0pt,
    left=8pt, right=8pt, top=8pt, bottom=6pt,
    borderline north={2pt}{0pt}{fachfarbe},
    before skip=6pt, after skip=6pt,
    fonttitle=\bfseries,
    #1
}

\newtcolorbox{formelkasten}{
    colback=white,
    colframe=fachfarbe,
    boxrule=1.5pt,
    arc=0pt,
    left=8pt, right=8pt, top=6pt, bottom=6pt,
    before skip=4pt, after skip=4pt
}

\newtcolorbox{bnebox}{
    colback=white,
    colframe=bnefarbe,
    boxrule=1.5pt,
    arc=0pt,
    left=8pt, right=8pt, top=6pt, bottom=6pt,
    before skip=4pt, after skip=4pt
}

% ═══════════════════════════════════════════════════════════════
% BEFEHLE
% ═══════════════════════════════════════════════════════════════

\newcommand{\luecke}[1]{\underline{\hspace{#1}}}
\newcommand{\punktebox}[1]{\hfill\small\textcolor{mittelgrau}{[\_\_/#1]}}

\newcommand{\aufgabe}[1]{%
    \needspace{4\baselineskip}%
    \vspace{10pt}%
    \noindent\textbf{\textcolor{fachfarbe}{Aufgabe #1}}%
    \vspace{4pt}%
    \nopagebreak[4]%
}

\newcommand{\operator}[1]{\textbf{\textcolor{fachfarbe}{#1}}}

\newlist{teil}{enumerate}{2}
\setlist[teil,1]{label=\alph*), leftmargin=1.8em, itemsep=3pt, parsep=0pt, topsep=3pt}

\newcommand{\antwortzeilen}[1]{%
    \par\vspace{4pt}%
    \foreach \n in {1,...,#1}{%
        \noindent\rule{\linewidth}{0.3pt}\vspace{6pt}%
    }%
}

\setlength{\parindent}{0pt}
\setlength{\parskip}{4pt}

% ═══════════════════════════════════════════════════════════════
% DOKUMENT
% ═══════════════════════════════════════════════════════════════

\begin{document}

% ═══════════════════════════════════════════════════════════════
% HEADER
% ═══════════════════════════════════════════════════════════════

\begin{center}
    {\Large\bfseries\textcolor{fachfarbe}{\thema}}\\[0.2em]
    {\small Arbeitsblatt | \fach{} Klasse \klasse{} | Stunde \stunde}
\end{center}
\vspace{-0.3em}
\hrule
\vspace{0.6em}

% ═══════════════════════════════════════════════════════════════
% TEIL 1: MERKKÄSTEN
% ═══════════════════════════════════════════════════════════════

\textbf{\large Teil 1: Merkkästen}

\begin{merkkasten}[title={\textbf{Kasten 1: Energieformen}}]
\begin{tabular}{@{}p{3.5cm}p{8cm}@{}}
\textbf{Lageenergie} $E_\text{pot}$ & Energie aufgrund der \luecke{2.5cm} eines Körpers \\[8pt]
\textbf{Bewegungsenergie} $E_\text{kin}$ & Energie aufgrund der \luecke{2.5cm} eines Körpers \\[8pt]
\textbf{Thermische Energie} $E_\text{therm}$ & Energie aufgrund der \luecke{2.5cm} der Teilchen \\[8pt]
\textbf{Elektrische Energie} $E_\text{el}$ & Energie durch bewegte \luecke{2.5cm} \\[8pt]
\textbf{Chemische Energie} $E_\text{chem}$ & In \luecke{2.5cm} gespeicherte Energie
\end{tabular}
\end{merkkasten}

\begin{merkkasten}[title={\textbf{Kasten 2: Energieerhaltungssatz}}]
\begin{formelkasten}
\centering
\Large $E_\text{ges} = E_1 + E_2 + E_3 + \ldots = \text{konstant}$
\end{formelkasten}

\vspace{4pt}
Energie kann nicht erzeugt oder \luecke{2.5cm} werden.

Sie wird nur von einer Form in eine andere \luecke{3cm}.

\vspace{6pt}
\textbf{Gültigkeitsbedingung:} Der Erhaltungssatz gilt nur in einem \luecke{3.5cm} System

(= keine Energie fließt von außen hinein oder nach außen ab).
\end{merkkasten}

\begin{merkkasten}[title={\textbf{Kasten 3: Energieflussdiagramm}}]
Ein Energieflussdiagramm zeigt, welche Energieformen in welche umgewandelt werden.

\vspace{6pt}
\textbf{Fadenpendel – drei Positionen:} \textit{Beschrifte die Energieformen an jeder Position!}

\begin{center}
\begin{tikzpicture}[scale=0.9]
    % Aufhängung
    \fill[gray!60] (0,3.2) rectangle (8,3.5);
    \draw[thick] (0,3.2) -- (8,3.2);

    % Position L (links) – Pendel ausgelenkt
    \begin{scope}[shift={(1.2,0)}]
        \draw[thick] (0,3.2) -- (-1.2,1.2);
        \fill[fachfarbe] (-1.2,1.2) circle (0.35);
        \node at (-1.2,0.4) {\footnotesize \textbf{L}};
        \node at (-1.2,-0.1) {\footnotesize $E_\text{pot}$: \luecke{1cm}};
        \node at (-1.2,-0.5) {\footnotesize $E_\text{kin}$: \luecke{1cm}};
    \end{scope}

    % Position M (Mitte) – Pendel unten
    \begin{scope}[shift={(4,0)}]
        \draw[thick] (0,3.2) -- (0,1);
        \fill[fachfarbe] (0,1) circle (0.35);
        \node at (0,0.2) {\footnotesize \textbf{M}};
        \node at (0,-0.3) {\footnotesize $E_\text{pot}$: \luecke{1cm}};
        \node at (0,-0.7) {\footnotesize $E_\text{kin}$: \luecke{1cm}};
    \end{scope}

    % Position R (rechts) – Pendel ausgelenkt
    \begin{scope}[shift={(6.8,0)}]
        \draw[thick] (0,3.2) -- (1.2,1.2);
        \fill[fachfarbe] (1.2,1.2) circle (0.35);
        \node at (1.2,0.4) {\footnotesize \textbf{R}};
        \node at (1.2,-0.1) {\footnotesize $E_\text{pot}$: \luecke{1cm}};
        \node at (1.2,-0.5) {\footnotesize $E_\text{kin}$: \luecke{1cm}};
    \end{scope}

    % Pfeile für Bewegung
    \draw[->, thick, gray] (0.5,2) -- (2.5,1.4);
    \draw[->, thick, gray] (5.5,1.4) -- (7.5,2);
\end{tikzpicture}
\end{center}

\vspace{4pt}
\textbf{Energiefluss (ideal):} \textit{Trage die Energieformen in die Kästchen ein!}

\begin{center}
\begin{tikzpicture}[
    box/.style={draw, thick, minimum width=2cm, minimum height=0.8cm, fill=white},
    arr/.style={->, >=Stealth, thick, line width=1.5pt}
]
    \node[box] (epot1) at (0,0) {\luecke{1.2cm}};
    \node[box] (ekin) at (3.5,0) {\luecke{1.2cm}};
    \node[box] (epot2) at (7,0) {\luecke{1.2cm}};

    \draw[arr] (epot1) -- node[above]{\small L→M} (ekin);
    \draw[arr] (ekin) -- node[above]{\small M→R} (epot2);
\end{tikzpicture}
\end{center}

\vspace{4pt}
\textbf{In der Realität:} Ein Teil der Energie wird durch \luecke{2.5cm} in \luecke{2.5cm} umgewandelt.

Deshalb schwingt das Pendel immer \luecke{2.5cm} und kommt schließlich zur Ruhe.
\end{merkkasten}

\begin{merkkasten}[title={\textbf{Kasten 4: Wirkungsgrad}}]
\begin{formelkasten}
\centering
\Large $\eta = \dfrac{E_\text{nutz}}{E_\text{zu}} = \dfrac{P_\text{nutz}}{P_\text{zu}}$
\end{formelkasten}

\vspace{4pt}
Der Wirkungsgrad $\eta$ (eta) gibt an, welcher Anteil der zugeführten Energie \luecke{2cm} wird.

\textit{Alternativ kann statt der Energie auch die Leistung $P$ eingesetzt werden.}

$\eta$ liegt immer zwischen \luecke{1cm} (0\%) und \luecke{1cm} (100\%).

\vspace{6pt}
\textbf{Beispiele:}
\begin{itemize}[leftmargin=1.5em, itemsep=2pt]
\item Elektromotor: $\eta \approx$ \luecke{1.5cm}
\item LED: $\eta \approx$ \luecke{1.5cm}
\item Glühbirne: $\eta \approx$ \luecke{1.5cm}
\end{itemize}

\vspace{4pt}
\textbf{Kettenrechnung:} $\eta_\text{ges} = \eta_1 \cdot \eta_2 \cdot \eta_3 \cdot \ldots$
\end{merkkasten}

% ═══════════════════════════════════════════════════════════════
% TEIL 2: AUFGABEN
% ═══════════════════════════════════════════════════════════════

\newpage

\textbf{\large Teil 2: Aufgaben}

\aufgabe{1 – Energieumwandlungen beschreiben}

($\star$) \operator{Ergänze} die fehlenden Energieformen in den Diagrammen. \punktebox{4}

\begin{teil}
\item \textbf{Ball fällt herunter (ideal):}

\vspace{4pt}
\begin{center}
\begin{tikzpicture}[
    box/.style={draw, thick, minimum width=2.2cm, minimum height=0.7cm},
    arr/.style={->, >=Stealth, thick, line width=1.5pt}
]
    \node[box] (a) at (0,0) {\luecke{1.5cm}};
    \node[box] (b) at (4.5,0) {\luecke{1.5cm}};
    \draw[arr] (a) -- (b);
    \node[below=0.2cm of a] {\small oben};
    \node[below=0.2cm of b] {\small unten};
\end{tikzpicture}
\end{center}

\item \textbf{Auto bremst ab:}

\vspace{4pt}
\begin{center}
\begin{tikzpicture}[
    box/.style={draw, thick, minimum width=2.2cm, minimum height=0.7cm},
    arr/.style={->, >=Stealth, thick, line width=1.5pt}
]
    \node[box] (a) at (0,0) {\luecke{1.5cm}};
    \node[box] (b) at (4.5,0) {\luecke{1.5cm}};
    \draw[arr] (a) -- (b);
    \node[below=0.2cm of a] {\small vorher};
    \node[below=0.2cm of b] {\small nachher};
\end{tikzpicture}
\end{center}
\end{teil}

\vspace{0.5em}

\aufgabe{2 – Lageenergie berechnen (Pendel)}

($\star\star$) Ein Fadenpendel hat eine Masse von $m = \SI{2}{kg}$ und wird $h = \SI{0{,}5}{m}$ angehoben. \punktebox{6}

\vspace{4pt}
\textbf{Gegeben:} $m = \SI{2}{kg}$, \quad $h = \SI{0{,}5}{m}$, \quad $g = \SI{10}{\frac{m}{s^2}}$

\begin{teil}
\item \operator{Berechne} die Lageenergie $E_\text{pot}$ des Pendels in der Startposition.

\vspace{4pt}
\textbf{Formel:} $E_\text{pot} = m \cdot g \cdot h$

\vspace{6pt}
\textbf{Rechnung:}

\vspace{2.5em}

\textbf{Ergebnis:} $E_\text{pot} = $ \luecke{2cm}

\vspace{4pt}

\item Das Pendel schwingt los. Im tiefsten Punkt (Höhe $h = 0$) ist $E_\text{pot} = 0$.

\operator{Bestimme} die Bewegungsenergie $E_\text{kin}$ im tiefsten Punkt (ideales Pendel).

\vspace{6pt}
\textbf{Begründung:}

\vspace{2em}

\textbf{Ergebnis:} $E_\text{kin} = $ \luecke{2cm}
\end{teil}

\vspace{0.5em}

\aufgabe{3 – Geschwindigkeit aus Energieerhaltung}

($\star\star$) Das Pendel aus Aufgabe 2 schwingt durch den tiefsten Punkt. \punktebox{6}

\begin{teil}
\item \operator{Berechne} die Geschwindigkeit $v$ im tiefsten Punkt.

\vspace{4pt}
\textbf{Hinweis:} Aus $E_\text{kin} = \frac{1}{2} m \cdot v^2$ folgt: $v = \sqrt{\frac{2 \cdot E_\text{kin}}{m}}$

\vspace{6pt}
\textbf{Rechnung:}

\vspace{3em}

\textbf{Ergebnis:} $v = $ \luecke{2cm}

\vspace{4pt}

\item \textbf{Rückwärts-Aufgabe:} Ein anderes Pendel ($m = \SI{0{,}5}{kg}$) erreicht im tiefsten Punkt $v = \SI{4}{\frac{m}{s}}$.

\operator{Berechne}, aus welcher Höhe $h$ es losgelassen wurde.

\vspace{6pt}
\textbf{Rechnung:}

\vspace{3em}

\textbf{Ergebnis:} $h = $ \luecke{2cm}
\end{teil}

\vspace{0.5em}

\aufgabe{4 – Wirkungsgrad berechnen}

($\star\star$) Eine LED-Lampe wird mit $P_\text{zu} = \SI{10}{W}$ betrieben und gibt $P_\text{Licht} = \SI{4}{W}$ als Licht ab. \punktebox{5}

\begin{teil}
\item \operator{Berechne} den Wirkungsgrad $\eta$ der LED.

\vspace{4pt}
\textbf{Formel:} $\eta = \dfrac{P_\text{nutz}}{P_\text{zu}}$

\vspace{6pt}
\textbf{Rechnung:}

\vspace{2em}

\textbf{Ergebnis:} $\eta = $ \luecke{1.5cm} $= $ \luecke{1.5cm} \%

\vspace{4pt}

\item \operator{Berechne}, wie viel Leistung als Wärme abgegeben wird.

\vspace{6pt}
\textbf{Ergebnis:} $P_\text{Wärme} = $ \luecke{2cm}
\end{teil}

\vspace{0.5em}

\aufgabe{5 – Fehlersuche}

($\star\star\star$) Florian behauptet: \punktebox{4}

\begin{center}
\fbox{\parbox{0.85\textwidth}{\centering
\textit{„Ein Stein (m = 2 kg) fällt aus h = 5 m Höhe. \\
Seine Geschwindigkeit beim Aufprall ist: \\
$E_\text{pot} = m \cdot g \cdot h = 2 \cdot 10 \cdot 5 = 100$ J \\
$E_\text{kin} = \frac{1}{2} m \cdot v^2$ \quad $\Rightarrow$ \quad $v = \sqrt{E_\text{kin}} = \sqrt{100} = \SI{10}{\frac{m}{s}}$"}
}}
\end{center}

\begin{teil}
\item \operator{Finde} den Fehler in Florians Rechnung und \operator{erkläre}, was er falsch gemacht hat.

\antwortzeilen{2}

\item \operator{Berechne} die richtige Geschwindigkeit.

\vspace{6pt}
\textbf{Ergebnis:} $v = $ \luecke{2cm}
\end{teil}

\vspace{0.5em}

\aufgabe{6 – Kettenaufgabe: Windkraftanlage}

($\star\star\star$) Eine Windkraftanlage nimmt $E_\text{Wind} = \SI{1000}{J}$ Windenergie auf. \punktebox{5}

Der Rotor hat $\eta_1 = 0{,}40$, der Generator $\eta_2 = 0{,}90$.

\begin{teil}
\item \operator{Berechne} den Gesamtwirkungsgrad $\eta_\text{ges}$.

\vspace{4pt}
\textbf{Rechnung:} $\eta_\text{ges} = \eta_1 \cdot \eta_2 = $ \luecke{3cm}

\vspace{4pt}

\item \operator{Berechne}, wie viel elektrische Energie $E_\text{el}$ erzeugt wird.

\vspace{6pt}
\textbf{Rechnung:}

\vspace{2em}

\textbf{Ergebnis:} $E_\text{el} = $ \luecke{2cm}
\end{teil}

% ═══════════════════════════════════════════════════════════════
% FOOTER
% ═══════════════════════════════════════════════════════════════

\vfill
\begin{center}
\small\textcolor{mittelgrau}{Gesamtpunkte: \_\_\_ / 30 BE}
\end{center}

\end{document}
